\section{Discussion}
\label{sec:discussion}

\paragraph{What the guarantee does and does not say}
Theorem~\ref{thm:gate} bounds the realised rate of the \emph{declared}
loss on the \emph{served} stream. It does not say that the declared loss
captures what the operator cares about; a contract on token-F1 is a
contract on token-F1. It does not promise anything about a single
answer, which is the nature of a rate. And it says nothing about the
queries that were declined -- if $40\%$ of traffic is handed to a human,
the guarantee covers the $60\%$ that was served, and the decline rate is
reported in every table precisely because it is the other half of the
picture. An operator reading Table~\ref{tab:main} should read cost, risk
and decline rate together; a low cost bought with a $60\%$ decline rate
is a different service from a low cost bought with $5\%$.

\paragraph{The auditor is the real assumption}
Every accounting scheme is only as good as its ledger entries.
Theorem~\ref{thm:gate} needs $\ell_t$, and $\ell_t$ comes from an
auditor -- an LLM judge, an NLI model, a human. If the auditor is
systematically biased, the ledger is systematically wrong and no amount
of bookkeeping repairs it; this is a strictly stronger dependence than a
calibration certificate has, since that one needs labels only during
calibration. We regard this as the honest cost of a deterministic
guarantee and have tried to price it rather than hide it:
Section~\ref{sec:res:audit} measures what partial labelling costs under
both accounting rules, and the audit line is inside every cost figure in
the paper. The practical reading is that a deterministic contract is
appropriate where auditing is already happening -- regulated settings,
services with human review queues -- and that elsewhere the
inverse-probability variant, with its probabilistic guarantee and its
measured $6$--$59\%$ breach rate, is the available compromise.

\paragraph{Why the ledger does not simply exhaust itself}
A natural worry is that the gate is a one-way ratchet: violations
accumulate, the ledger empties, and the system declines forever. It does
not, because declining earns allowance at rate $\alpha$. If the deployed
mixture has risk $r > \alpha$, the steady-state decline rate settles at
$1 - \alpha/r$, which is exactly the minimum any feasible policy must
have (Proposition~\ref{prop:feasible}); the controller discovers it
without being told $r$. The measured decline rates track this
prediction: on QAMPARI at $\alpha = 0.35$ with a risk floor of $0.591$,
the bound is $40.8\%$ and the executor declines $51.9\%$, the gap being
the cheaper, riskier plans it chooses to run in place of the safest one.

\paragraph{When this is the wrong design}
Three regimes favour the incumbent. If labels are unavailable after
deployment at any price, the ledger cannot be maintained and a
calibration certificate is the only option. If the stream is short --
shorter than the few hundred queries the estimates need -- the
controller is still in its transient and a certificate calibrated
offline may be cheaper; our ASQA results, on the shortest stream, show
the smallest margin and the largest sensitivity to $\kappa$. And if
declining is genuinely impossible, so that $\mathcal{A} = \mathcal{P}$,
the gate has no safe action to fall back on and
Theorem~\ref{thm:gate} does not apply; the best available then is the
probabilistic guarantee, and the contract is capped at $\min_p r_p$
whatever one does.

\paragraph{Multi-dimensional contracts}
A service typically owes several rates at once -- quality, citation
support, freshness, latency. The construction extends directly: run one
ledger per dimension and open the gate only when all of them can absorb
a violation, which yields $V_t^{(j)} \le \alpha_j t$ simultaneously for
every $j$ with no union correction, since each bound is deterministic.
The optimisation layer becomes an LP with $d$ risk constraints, whose
basic solutions have at most $d+1$ nonzero weights. This is a strict
simplification over the intersection--union testing a certified
formulation needs, and it is the clearest illustration of what the
reframing buys: the multiplicity problems that dominate the statistical
design disappear when the guarantee stops being statistical. We report
the single-dimension case here because our execution matrices carry
per-query costs for one auditable loss per benchmark; the multi-ledger
extension is implemented and left for evaluation on a workload with
several audited dimensions.

\paragraph{Threats to validity}
Costs are list prices at a point in time and vendor pricing moves; we
report the metering method and release the token counts so any table can
be repriced. GPU-second pricing for local models depends on an assumed
rental rate. The streams are permutations of finite pooled populations
rather than fresh traffic, so the long-horizon results resample the same
queries; this is why we report regret against horizon rather than
claiming a steady state. Finally, the plan grids cover two of the four
benchmarks, since executing $36$ and $24$ plans end to end on all four
was beyond our budget; the four-benchmark results use the
$69$-candidate space over the ladder. Two grids are enough to show that
the effect of enlarging the space has no fixed sign, but not enough to
predict its sign from properties of a task, which we regard as the more
useful open question this raises.
